\chapter{Perfil de Egreso}

\section{Justificación}

Durante las dos primeras décadas del siglo XXI, el desarrollo de la matemática está en constante aumento gracias a la digitalización de la mayor parte de los objetos usados en el siglo pasado (fotografía, música, medicina, documentos, etc.) y su transferencia.  El célebre matemático Cedric Villani, a quien se le otorgó la medalla Fields, en el año 2010, mencionó en una conferencia en la Facultad de Ciencias de la Universidad Nacional Autónoma de México, que la cantidad de estudiantes de matemática ha aumentado más del 20\% en los últimos 20 años, y esto se debe principalmente a la digitalización del mundo. 

    La Organización de las Naciones Unidas para la Educación, la Ciencia y la Cultura (UNESCO,por sus siglas en inglés) menciona en [2], el importante desarrollo de la matemática a través del mundo, y, sobre todo, en los países en vías de desarrollo:
    
    “Through international cooperation and partnerships, UNESCO organizes activities, especially in developing countries, to promote capacity-building for research and advanced training in mathematics and mathematics education, and in general, to enhance public understanding and appreciation of the importance of mathematics in society and daily life.”

y

    “The UNESCO International Mathematics Exhibition: Experiencing Mathematics continues to travel around the world since July 2004. The exhibition was conceptualized and designed to promote public appreciation of the importance of mathematics to show that mathematics is indispensable for daily living, from basic economic activities to the operation of train stations and airports, as well as to demonstrate that mathematics is fun and interesting. Experiencing Mathematics features posters, manipulative models and interactive devices.”
    
Debido a la importancia de las matemáticas y su interacción con otras áreas, es que la cuadragésima Conferencia General de la UNESCO, en noviembre de 2019, proclamó el 14 de marzo de cada año como el Día Internacional de las Matemáticas, y la siguiente cita figura en su página web, cf. [3]: 

    “Una mayor conciencia mundial y un fortalecimiento de la enseñanza de las ciencias matemáticas son esenciales para hacer frente a desafíos que se plantean en ámbitos como la inteligencia artificial, el cambio climático, la energía y el desarrollo sostenible, y para mejorar la calidad de vida en el mundo desarrollado y en el mundo en desarrollo”. 

    Además, ahí se menciona adicionalmente: 
    “El Día Internacional de las Matemáticas tiene por objeto destacar el papel fundamental que desempeñan las ciencias matemáticas en el logro de los Objetivos de Desarrollo Sostenible de las Naciones Unidas y en el fortalecimiento de las dos prioridades de la UNESCO: África y la igualdad de género. El Día nos invita a celebrar la alegría de las matemáticas y la plétora de vocaciones que ofrece a niñas y niños a través de actividades festivas y diversas en todo el mundo.”
    
 Considerando lo expuesto por la UNESCO en la cuadragésima conferencia, la temática de este año 2020 es: Las matemáticas están en todas partes. En [3] se lee:
 
    ‘’El tema de este año … transmite que las matemáticas son parte del patrimonio cultural de la humanidad: son tan esenciales para nuestras tecnologías como una herramienta para el desarrollo. Las numerosas actividades que tienen lugar en todo el mundo para celebrar el Día son testimonio de la entusiasta acogida que ha tenido. El Día Internacional de las Matemáticas está dirigido por el Programa Internacional de Ciencias Básicas de la UNESCO (PICF) y la Unión Matemática Internacional, con el apoyo de numerosas organizaciones internacionales y regionales.”
    
Y la importancia de la matemática en varias áreas y actividades la humanidad se ven en la ciencia y la tecnología, cuyos ejemplos son variados, según [3]:

\begin{itemize}

  \color{blue}  \item \color{black} El éxito de los buscadores de internet viene de su brillante algoritmo matemático
    
   \color{blue} \item \color{black} La criptografía para comunicaciones seguras se basa en la teoría de números.   
    
    \color{blue} \item \color{black} Los dispositivos de imagen médica, como escáneres de tomografía computarizada (CT-scan) o de imagen por resonancia magnética (MRI), recogen datos numéricos y un algoritmo matemático construye una imagen a partir de ellos.   

    \color{blue} \item \color{black} La inteligencia artificial y el aprendizaje automático (machine learning) están transformando el mundo; por ejemplo, la visión artificial, traducción automática, vehículos autónomos, etc.
    
    \color{blue} \item \color{black} Descodificar el genoma humano es un triunfo de las matemáticas, la estadística y la informática.
    
    \color{blue} \item \color{black} Las matemáticas se usaron para crear la primera imagen de un agujero negro.
    
    \end{itemize}
    
    Como herramienta de organización en la sociedad, se reconoce el papel de la matemática,
tal y como se presenta en [3]:

\begin{itemize}
    
     \color{blue} \item \color{black} Las matemáticas se usan para optimizar el transporte y las redes de comunicaciones
    
    \color{blue} \item \color{black} Las matemáticas ayudan a comprender y controlar la propagación de epidemias. 
    
    \color{blue} \item \color{black} La estadística y la optimización se usan en la planificación y gestión eficiente de sistemas de salud pública, economía, y sociedad. 
    
    \color{blue} \item \color{black} Las matemáticas ayudan a diseñar sistemas electorales que mejor representen la voluntad popular.   
    
    \color{blue} \item \color{black} Las matemáticas ayudan a comprender los riesgos naturales (inundaciones,  terremotos,
huracanes) y a prepararse con anticipación para evitar desastres

\end{itemize}

La organización de las Naciones Unidas, a través de la UNESCO, menciona que las matemáticas son esenciales para alcanzar los Objetivos de Desarrollo Sostenible de la ONU, y en alusión a esto, mencionan en [3], los siguientes ejemplos: 

\begin{itemize}
    
     \color{blue} \item \color{black} Las matemáticas son una herramienta para el desarrollo. Citando a Nelson Mandela (Junio 1990), “La educación es el arma más poderosa que se puede usar para cambiar el mundo”. Y las matemáticas son una parte esencial de la educación, que es necesaria para tener mejores trabajos. 
    
    \color{blue} \item \color{black} Las matemáticas se usan para modelar cambios globales y sus consecuencias en la biodiversidad.
    
    \color{blue} \item \color{black} Las matemáticas se usan para modelar cambios globales y sus consecuencias en la biodiversidad.
    
    \color{blue} \item \color{black} Técnicas de optimización y análisis de datos son necesarias en la transición a un uso sostenible de los recursos naturales.
    
    \color{blue} \item \color{black} La inteligencia artificial se utiliza para extraer datos de imágenes de satélite y dibujar mapas de áreas urbanas, industriales, agrícolas y forestales donde no se encuentran datos por métodos tradicionales.   

    \color{blue} \item \color{black} La educación en matemáticas ofrece a las chicas y a las mujeres un futuro mejor.   

    \color{blue} \item \color{black} La cultura científica y matemática ayuda a cada ciudadano a comprender mejor los desafíos que le rodean. 
    
     \color{blue} \item \color{black} Finalmente, en esta mención, [3], de lo que expresa la UNESCO en relación con las matemáticas, se manifiesta que están en todas partes y en todo lo que hacemos: 
    
     \color{blue} \item \color{black} Las matemáticas inspiran a artistas y músicos: perspectiva, simetría, mosaicos, fractales, curvas, superficies y formas geométricas; patrones, escalas y sonidos en música. 
      
     \color{blue} \item \color{black} Las matemáticas son útiles en juegos de estrategia, desde el backgammon o el ajedrez hasta resolver un cubo de Rubik o jugar al Awale.   


    \color{blue} \item \color{black} Las matemáticas son útiles para preparar presupuestos.
     
      
    \color{blue} \item \color{black} Prácticamente todo el mundo usa algunos conceptos matemáticos: el constructor, el granjero, el comerciante, el artesano, el atleta , …
    
    \color{blue} \item \color{black} Las matemáticas están detrás de las técnicas de geolocalización, desde la navegación con  las estrellas hasta el GPS.   
    
    \color{blue} \item \color{black} Las matemáticas están detrás del software de nuestros teléfonos.   
     
      
    \color{blue} \item \color{black} Las matemáticas hacen películas de animación realistas. 
    
    \color{blue} \item \color{black} ¿Visitaremos Marte algún día? Sin matemáticas nunca será posible.
    \end{itemize}

Todo esto nos muestra que la matemática ha sido usada a lo largo de la historia de la civilización en diferentes circunstancias. Cabe preguntarnos algunas de las siguientes proezas, tal y como se lo preguntan en [4]: 

    « ¿cómo se construyeron las pirámides en Egipto? ¿Cómo se suspenden los puentes? ¿Cómo fue el aterrizaje en la Luna? » Pero, más allá de estos hitos que han marcado a la historia y al desarrollo de las sociedades, resaltamos un interés que no ha sido mencionado: las habilidades que desarrolla un estudiante que incursiona en cualquier carrera de matemática en cualquier parte del mundo.  

La matemática es excelente para el cerebro, puesto que desarrolla importantes conexiones neurológicas para procesar la información, y, entonces, no es sorprendente que estos cerebros logren habilidades analíticas que son fuertemente buscadas y deseadas por los empleadores. En este sentido, los investigadores conducidos por la Dra. Tanya Evans, de la Universidad de Stanford, mencionan que los niños que saben matemáticas son capaces de hacer funcionar ciertas regiones del cerebro mejor y tienen materia gris más voluminosa en estas regiones.  De esta forma se evidencia
en [4] la siguiente cita:

    “Children who know math are able to recruit certain brain regions more reliably, and have greater gray matter volume in those regions, than those who perform more poorly in math. The brain regions involved in higher math skills in high-performing children were associated with various cognitive tasks involving visual attention and decision-making.” 

La matemática lleva a solucionar problemas de la vida real. Ya hemos visto, anteriormente, muchos ejemplos de aplicaciones de la matemática a la vida cotidiana, lo que es interesante es que el matemático, en su estructura mental, enfrentará cada problema de su cotidianidad intentando aplicar el rigor para solventar los problemas que se presentan. Y en algunos casos, podría dar soluciones novedosas. Lo que también es cierto es que los matemáticos tienen un entrenamiento diario en resolución de problemas, y se desarrolla una capacidad no solo para solventar el problema, sino, más bien, para plantearlo e inicialmente para comprenderlo, lo cual no necesariamente cualquiera intenta desde el principio. Definitivamente, los estudiantes de matemática desarrollan un mejor sistema para resolver problemas a lo interno y externo de las matemáticas.

La matemática ayuda en casi todas las demás carreras universitarias.  La complejidad y diversidad de las matemáticas es la razón principal, por la cual, otras disciplinas se benefician de este conocimiento y su forma de argumentar y razonar.  En este sentido, se puede ver en [4] que la Asociación Americana de Matemática (AMS, por sus siglas en inglés) menciona: 

    “Mathematical research and education are at the heart of some careers, while other careers utilize mathematics and its applications to build and enhance important work in the sciences, business, finance, manufacturing, communications, and engineering”

Las carreras que más se benefician de la matemática son principalmente las ingenierías, la computación, la física, la estadística, las ciencias actuariales, las finanzas, la economía y la administración de negocios. Sin embargo, existen otras carreras que se pueden beneficiar colateralmente de las matemáticas, como son las letras y las artes, en un sentido más general; como ejemplo podemos mencionar a Jorge Luis Borges con su extensa literatura y la proporción áurea en las artes.  En este sentido, la matemática usualmente es referida como “la reina de las ciencias”.

La matemática ayuda a entender mejor el mundo.  Cuando se entienden las matemáticas se llega a comprender esencialmente fenómenos naturales. Por ejemplo, la teoría de relatividad de Einstein es la curvatura del espacio cuatro dimensional formado por el tiempo y el espacio tridimensional en el que vivimos.  Pero los conceptos concretos se encuentran en la geometría diferencial.  De ahí que el estudio y comprensión de esa geometría son esenciales para el entendimiento de la teoría de la relatividad. 
La matemática es un lenguaje universal.  Independientemente del lenguaje materno que se hable, el matemático podrá leer y entender la matemática en reuniones científicas o en medios escritos por sí mismo, sin necesariamente hablar la lengua en la cual se escribe. Esto refuerza el hecho de que la matemática está constituida de objetos, relaciones y descripciones universales que no dependen tanto de la cultura e idiosincrasia.  

La matemática y las carreras que se relacionan con ella han sido una de las mejores clasificadas en “CareerCast.com\'s 2016 rankings of the 200 best jobs”. Por ejemplo, mencionamos el número en donde quedaron clasificadas las siguientes carreras: en primer lugar, analista de datos; y en segundo, ser estadístico. En sexto lugar, ser matemático y en décimo puesto, ser un actuario. Y carreras que tienen un alto contenido matemático quedaron respectivamente en las siguientes posiciones: Ciencias de la computación (\# 7), Ingenieros biomecánicas ( \# 14), Meteorología (\# 16), Físicos (\# 21), Economistas (\# 24), Analistas financieros (\# 27), Ingenieros ambientales (\# 37), Ingenieros civiles (\# 38), Ingenieros mecánicos (\# 58), e Ingenieros aeroespaciales (\# 59). 

Además, la Oficina de Trabajo Estadístico (The Bureau of Labor Statistics) proyecta que los empleos en ocupaciones matemáticas crecerá un 28\%, entre los años 2016 al 2026, crecimiento más rápido que el promedio de crecimiento de todas los demás empleos. Esto se debe a la relación que tienen estos trabajos con el “Big data”. 

Por todas estas razones y por la constitución humanista de la Universidad de Costa Rica que se asienta en su estatuto orgánico, es que ha sido, es y será relevante para la institución y para la nación que se impartan carreras de matemática con diferentes énfasis.



\subsection{Perfil de salida}

En la presente sección se presenta una síntesis del perfil de egreso elaborado en el periodo 2018–2020 por la Comisión de Malla Curricular de la carrera Bach. y Lic. en Matemática, como parte del proceso de transformación curricular que impulsa la Escuela de Matemática. El contenido completo del perfil aprobado se incluye en los anexos de este documento.



\begin{table}[H]
\centering
\renewcommand{\arraystretch}{1.3}
\begin{tabular}{|c|p{9cm}|}
\hline
\cellcolor[HTML]{00C0F3}\textbf{Código(s)} & \cellcolor[HTML]{00C0F3}\textbf{Enunciado de la dimensión declarativa} \\
\hline
CD01 & Conoce y comprende los fundamentos matemáticos —incluyendo la lógica, la teoría de números y la geometría euclídea— y los principios formales que sustentan el razonamiento y la argumentación rigurosa. \\
\hline
CD02 & Conoce el uso de software especializado, los fundamentos de la programación y la computación científica, así como las técnicas básicas de análisis numérico para formular, explorar y resolver problemas matemáticos mediante enfoques computacionales. \\
\hline
CD03 & Conoce los conceptos, métodos y aplicaciones básicos del cálculo y del álgebra lineal como fundamentos de la formación matemática. \\
\hline
CD04 & Conoce los conceptos y técnicas de las ecuaciones diferenciales, la teoría de la medida, la probabilidad y los procesos estocásticos, así como del análisis complejo y funcional. \\
\hline
CD05 & Conoce los fundamentos y resultados básicos de la geometría y la topología. \\
\hline
CD06 & Conoce las estructuras y resultados fundamentales del álgebra abstracta. \\
\hline
CD07 & Conoce los principios de la matemática aplicada, por ejemplo por medio de la programación, la modelación, computación científica o el análisis de datos, para resolver problemas y tomar decisiones en ámbitos como la industria, la banca y otros sectores.
\\
\hline
\end{tabular}
\caption{Tabla nueva de conocimientos declarativos.}
\label{tab:declarativa-egreso-integrada}
\end{table}

%\newpage

\begin{table}[H]
\centering
\renewcommand{\arraystretch}{1.3}
\begin{tabular}{|c|p{9cm}|}
\hline
\cellcolor[HTML]{00C0F3}\textbf{Código(s)}  & \cellcolor[HTML]{00C0F3}\textbf{Enunciado de la dimensión procedimental} \\
\hline
CP01 & Aplica pensamiento operacional, abstracto y formal para analizar, resolver y sintetizar problemas matemáticos. \\
\hline
CP02 & Formula conjeturas, resuelve problemas, interpreta resultados y se comunica de manera coherente en lenguaje matemático formal. \\
\hline
CP03 & Utiliza tecnologías y software científico para explorar, validar y resolver problemas matemáticos.\\
\hline
CP04 & Participa activamente en discusiones científicas y divulga el conocimiento matemático de forma accesible a diferentes audiencias.\\
\hline
\end{tabular}
\caption{Tabla nueva de conocimientos procedimentales.}
\label{tab:procedimental-egreso-integrada}
\end{table}


\newpage

\begin{table}[H]
\centering
\renewcommand{\arraystretch}{1.3}
\begin{tabular}{|c|p{9cm}|}
\hline
\cellcolor[HTML]{00C0F3}\textbf{Código(s)} & \cellcolor[HTML]{00C0F3}\textbf{Enunciado de la dimensión actitudinal} \\
\hline
CA01 & Manifiesta respeto por la diversidad humana, matemática, científica y cultural, fomentando un ambiente inclusivo y de valoración mutua. \\
\hline
CA02 & Se comunica con claridad, colabora en equipos académicos y profesionales, y mantiene una disposición al autoaprendizaje y al liderazgo.\\
\hline
CA03 & Demuestra conciencia social y participación ciudadana, con una visión sociocultural e integral de la sociedad en la que la ciencia y la matemática ocupan un lugar relevante.\\
\hline
CA04 & Ejerce pensamiento crítico, actúa con integridad ética y evalúa con rigor las prácticas científicas y académicas.\\
\hline
CA05 & Utiliza su conocimiento matemático para resolver problemas académicos y profesionales, adaptándolo a distintos contextos y necesidades.\\
\hline
CA06 & Mantiene un interés permanente por los desarrollos de la investigación científica y matemática, valorando la actualización como parte de su formación profesional.\\
\hline
\end{tabular}
\caption{Tabla nueva de conocimientos actitudinales.}
\label{tab:actitudinal-egreso-integrada}
\end{table}





